\documentclass[12pt]{article}

% --- Formatting per candidacy requirements -------------------------------
% Proposition limit: 15 double-spaced pages, excluding front matter
% and references.
% This is the IN-FIELD proposition.
% -------------------------------------------------------------------------


\usepackage[letterpaper,margin=1in]{geometry}
\usepackage{setspace}
\doublespacing

\usepackage[T1]{fontenc}
\usepackage{lmodern}
\usepackage{microtype}

\usepackage{amsmath,amssymb,amsthm,mathtools}
\usepackage{physics}
\usepackage{graphicx}
\usepackage{booktabs}
\usepackage{siunitx}
\usepackage{xcolor}
\usepackage[hidelinks]{hyperref}
\usepackage[capitalize,noabbrev]{cleveref}

\usepackage[
  backend=biber,
  style=numeric-comp,
  sorting=none,
  giveninits=true,
  maxbibnames=10,
  url=false
]{biblatex}
\addbibresource{references.bib}

\graphicspath{{figures/}}

\title{Proposition 1 (In-Field) \\[0.25em]
       \large Tensor-Network Time Evolution for Spin Dynamics in
       Quantum Chemistry}
\author{Xiangzhi (Sam) Ouyang}
\date{\today}

\begin{document}

% --- Front matter (excluded from page count) -----------------------------
\pagenumbering{gobble}
\begin{titlepage}
  \centering
  \vspace*{2cm}
  {\Large \bfseries Candidacy Proposition 1\par}
  \vspace{0.5cm}
  {\large \textit{In-field}\par}
  \vspace{1.5cm}
  {\LARGE Tensor-Network Time Evolution for Spin Dynamics
          in Quantum Chemistry\par}
  \vspace{2cm}
  {\large Xiao-Yu Ouyang \par}
  \vspace{0.5cm}
  {\normalsize \texttt{xouyang@caltech.edu} \par}
  \vspace{1.5cm}
  {\normalsize Division of Chemistry and Chemical Engineering \\
               California Institute of Technology \par}
  \vfill
  {\normalsize Submitted \today \par}
\end{titlepage}

\pagenumbering{roman}
\begin{singlespace}
\tableofcontents
\end{singlespace}
\newpage

\begin{abstract}
\noindent
\textit{[One paragraph: central claim, target system, method,
expected result, why it is original to me and not a next-step
extension of group work.]}
\end{abstract}

% --- Body (counts toward the 15-page limit) ------------------------------
\pagenumbering{arabic}


\section{Central Claim}
This proposition develops graph-adapted tensor-network and
message-passing methods for finite-time dynamics of isolated spin
Hamiltonians on molecular coupling graphs.  The goal is to compute
experimentally relevant observables, such as magnetic-resonance
correlation functions, without forcing the molecular graph into a
one-dimensional chain ordering.

\section{Background and Motivation}

Many spectroscopic and chemical observables are determined not only by
static electronic structure, but by the coherent dynamics of effective
spin degrees of freedom.  Quantum chemistry can provide the parameters
of a spin Hamiltonian, including scalar $J$ couplings, hyperfine
couplings, dipolar couplings, exchange interactions, $g$ tensors, and
zero-field splittings.  The subsequent task is to propagate this spin
Hamiltonian and compute finite-time observables.  In this proposal I
focus on the isolated, coherent part of this problem,
\begin{equation}
  H_G =
  \sum_{i\in V} \mathbf h_i\cdot \mathbf S_i
  +
  \sum_{(i,j)\in E}
  \mathbf S_i\cdot \mathbf J_{ij}\cdot \mathbf S_j ,
  \label{eq:graph-spin-hamiltonian}
\end{equation}
where $G=(V,E)$ is a molecular spin-coupling graph.  Relaxation,
pumping, recombination, and other open-system effects are important in
many experiments, but they are not the central target of this
proposition; they can be treated as phenomenological broadening or as
future extensions.

The primary experimental motivation is coherent magnetic-resonance
spectroscopy, especially liquid-state and zero-field NMR.  In zero-field
NMR, the relevant Hamiltonian is often dominated by scalar coupling
terms, $2\pi\sum_{(i,j)}J_{ij}\mathbf I_i\cdot \mathbf I_j$, and the
measured signal is a real-time correlation function of magnetization
operators.  Although the initial density matrix is mixed, the high
temperature signal can be written in terms of a simple deviation density
operator, for example $\Delta\rho(0)\propto \sum_i\gamma_i I_i^z$, and a
closed-system evolution,
\begin{equation}
  s(t) =
  \operatorname{Tr}\!\left[
    \Delta\rho(0) U^\dagger(t) O U(t)
  \right].
  \label{eq:observable-contraction}
\end{equation}
The same formal structure appears in controlled benchmark spin graphs,
where one asks how local spin observables, correlation functions, or
operator fronts evolve on chains, trees, and sparse molecular-like
graphs.  A secondary motivation comes from radical-pair spin chemistry:
the coherent component is a two-electron spin system coupled to local
nuclear spin baths, with singlet-triplet dynamics generated by Zeeman,
hyperfine, exchange, and dipolar terms \cite{Fay2020RadicalPair,
Antill2024RadicalPy}.  In this proposal I use this as a second example
of chemically relevant isolated spin dynamics, while leaving
spin-selective recombination kinetics outside the core scope.

Direct simulation of these systems is limited by the exponential growth
of Hilbert or Liouville space.  Exact diagonalization, Krylov
propagation, and Chebyshev methods are reliable for small systems, but a
general $N$-spin Hilbert space scales as $\prod_i(2S_i+1)$, and a density
operator representation scales as its square.  Magnetic-resonance
simulation packages such as Spinach address this by using symmetry,
sparse Liouville-space representations, and state-space restriction; the
Spinach work reports routine liquid-state simulations for 40 or more
spins in favorable cases \cite{Hogben2011Spinach}.  These methods are
powerful, but their truncations are usually based on operator rank,
coherence order, graph distance, or other physically motivated
restrictions rather than on a tensor-network environment optimized for a
specific finite-time observable.

Tensor-network methods provide an alternative language for controlled
compression.  Matrix product states and density-matrix renormalization
group methods are mature for one-dimensional systems
\cite{Schollwock2011DMRG}; real-time variants include TEBD and the
time-dependent variational principle \cite{Haegeman2011TDVP}.  To apply
these methods to a molecular spin graph, however, one must choose a
one-dimensional ordering of the spins.  A Hamiltonian that is local on
the molecular graph can then become long-ranged on the MPS chain, and
the required bond dimension is controlled by how much graph correlation
or entanglement crosses each artificial chain cut.  This ordering
dependence is familiar in quantum-chemistry DMRG and is expected to be
at least as important for branched or sparse molecular spin graphs.

Tree tensor network states provide a controlled compromise between a
one-dimensional MPS and a fully loopy graph tensor network.  A tree
retains efficient contraction and DMRG-like optimization, while allowing
a hierarchical geometry that can better match branched molecular
structure \cite{Shi2006Classical,nakataniEfficientTreeTensor2013}.
Prior tree tensor network work in quantum chemistry treated fermionic
active-space electronic Hamiltonians on orbital tensor networks; in
contrast, the present proposal targets effective spin Hamiltonians on
molecular coupling graphs.  This avoids fermionic sign structure and
turns the central topology problem into the choice of a tree or
low-treewidth structure adapted to the spin-coupling graph and to the
observable being computed.

The topology is not determined by the Hamiltonian alone.  The weighted
coupling graph, with weights such as $\|\mathbf J_{ij}\|$, provides a
physically motivated initial tree: strongly coupled spins should usually
be close in the tensor network, and high-degree tensors should be
avoided for computational reasons.  The optimal compression, however, is
set by the entanglement or message spectra generated by the target
dynamics.  For the observable in \cref{eq:observable-contraction}, the
natural tensor network is a closed spacetime contraction.  In one spatial
dimension it has the geometry of a spacetime cylinder; for a molecular
spin graph, the spatial direction is replaced by the graph $G$.  This
suggests a graph-aware message-passing approach in which the contraction
is organized along a tree or sparse graph and intermediate messages are
compressed using the environment relevant to the final observable.  Such
an approach connects the practical goal of spin-dynamics simulation to
environment-dependent tensor-network compression and belief-propagation
ideas developed for real-time tensor-network contractions
\cite{Park2025IFBP}.










% --- References (excluded from page count) -------------------------------
\begin{singlespace}
\printbibliography
\end{singlespace}

\end{document}
